% Recuperación autocontenida: sector angular de árbol y reconstrucción Higgs--vacío.
% Propietarios y alcance: higgs_procedencia.md. Incorporar después del operador constitutivo.
\section{Reconstrucción conjunta del calibre, Higgs y la respuesta constitutiva}
\label{amp-higgs-seccion}

La respuesta del vacío y el sector electrodébil admiten lecturas complementarias
del transporte angular ya construido. La primera conserva los dos canales
marcados; el determinante conserva su contracción común; la componente angular
de Higgs conserva además la orientación que desaparece en ese determinante.
La composición de estas lecturas permite reconstruir su antecedente espectral
y expresar la impedancia mediante el autoacoplamiento escalar y los acoplamientos
de calibre, dentro de una misma realización de árbol.

\subsection{Procedencia angular y firmas de las rutas}

Se conserva la cadena APP--TRIT--TPK: la prolongación superviviente determina
la ruta, ésta produce su registro orientado y el lector del registro expresa
la firma angular. Las coordenadas $\pi=\pi_{\rm HMT}$ y
$\alpha=\alpha_{\rm HMT}$ son salidas previas de la misma generación;
$A$ y $C^*$ son sus representaciones angulares posteriores. En este bloque se
usan radianes,
\begin{equation}
 x=\frac{\pi A_{\deg}}{180},\qquad
 y=\frac{\pi C^*_{\deg}}{180},\qquad
 A_{\deg}=1000\alpha,\qquad 0<y<x.
 \label{amp-higgs-angulos}
\end{equation}
Por tanto, $x=50\pi\alpha/9$. Los canales y las marcas de hoja son los
de la sección~\ref{iii:sec:hojas}:
\begin{equation}
 \begin{aligned}
 q_+&=e^{-x-y},& q_-&=e^{-x+y},& D&=q_+q_-=e^{-2x},\\
 P_\pm&=\frac{I\pm\mathcal R}{2},&
 T&=q_+P_++q_-P_-.
 \end{aligned}
 \label{amp-higgs-transporte}
\end{equation}
Aquí $\mathcal R$ es la involución autoadjunta marcada y cada proyector
tiene rango uno en la realización bicapa básica.

El lector angular de duración--perduración es
\begin{equation}
 \chi_{\rm dur}(n,s)=\exp\!\left(nx+\frac{s}{6}y\right).
 \label{amp-higgs-caracter}
\end{equation}
Para conservar la procedencia finita de los coeficientes que se utilizarán,
reunimos las huellas nominales de las tres rutas del calendario de longitud
$108$. Las columnas de desplazamientos cardinales se leen como $(N,E,S,O)$;
la pareja siguiente conserva la hoja y el residuo expresados por esa ruta.
La tabla conserva, para cada ruta identificada, la huella y la firma
angular expresadas por la cadena anterior.
\begin{center}
\small
\begin{tabular}{c c c c c}
\hline
Sector & Ruta y eje & $(N,E,S,O)$ & Hoja/residuo & $(n,s)$\\
\hline
$W$ & $15/37$, E--O & $(23,44,11,30)$ & $(41,8)$ & $(94,-1)$\\
$Z$ & $20/23$, E--O & $(33,45,10,20)$ & $(57,7)$ & $(95,-1)$\\
$H$ & $24/29$, N--S & $(40,34,22,12)$ & $(47,10)$ & $(97,9)$\\
\hline
\end{tabular}
\end{center}
Las tres huellas tipadas conservan el identificador de ruta, la orientación
y el registro de procedencia de las firmas empleadas. La evaluación angular
comienza después de la extracción desde el registro enriquecido; la tabla
reúne sus salidas para componer los caracteres en esta realización.
Las marcas $W_\pm$ que siguen designan coordenadas de enrollamiento,
distintas del nombre de la especie $W$:
\begin{equation}
 \begin{aligned}
 W_\pm&=6n\pm s, &
 n&=\frac{W_++W_-}{12},& s&=\frac{W_+-W_-}{2},\\
 (n,s)_W&=(94,-1), & (W_+,W_-)_W&=(563,565),\\
 (n,s)_Z&=(95,-1), & (W_+,W_-)_Z&=(569,571),\\
 (n,s)_H&=(97,9), & (W_+,W_-)_H&=(591,573).
 \end{aligned}
 \label{amp-higgs-firmas}
\end{equation}
La matriz de la primera aplicación es
$\bigl(\begin{smallmatrix}6&1\\6&-1\end{smallmatrix}\bigr)$, con determinante
$-12$ e imagen $\{(u,v)\in\mathbb Z^2:u+v\in12\mathbb Z\}$.
En efecto, las fórmulas inversas anteriores son enteras precisamente en
esa imagen. El carácter conserva las dos orientaciones mediante
\begin{equation}
 \chi_{\rm dur}(n,s)
 =\left(q_+^{-W_+}q_-^{-W_-}\right)^{1/12}.
 \label{amp-higgs-caracter-canales}
\end{equation}
La prueba se obtiene tomando logaritmos: el numerador es
$W_+(x+y)+W_-(x-y)=12nx+2sy$.
Así quedan identificados la raíz duodécima y el divisor $6$ del lector.

\begin{proposition}[Cocientes angulares determinados por las firmas]
En una sección de escala común, con los factores de refinamiento compartidos
que permiten su cancelación en los cocientes, se cumple
\begin{equation}
 \frac{m_W^{\angle}}{m_Z^{\angle}}=e^{-x},\qquad
 \frac{m_H^{\angle}}{m_W^{\angle}}
 =e^{3x+5y/3},\qquad
 \left(\frac{m_H^{\angle}}{m_W^{\angle}}\right)^2
 =e^{6x+10y/3}.
 \label{amp-higgs-cocientes}
\end{equation}
\end{proposition}
\begin{proof}
Dividir los caracteres de $W$ y $Z$ resta sus firmas: $(94,-1)-(95,-1)=(-1,0)$.
Para $H/W$, la diferencia es $(97,9)-(94,-1)=(3,10)$.
La ecuación~\eqref{amp-higgs-caracter} da $3x+(10/6)y$; el cuadrado
multiplica ese exponente por dos. La escala común y los refinamientos
compartidos se cancelan antes de estas operaciones.
\end{proof}
Estos cocientes conservan el identificador de la coordenatización angular. La aplicación
a masas con refinamientos sectoriales distintos conserva sus factores
adicionales; la condición de cancelación forma parte del enunciado.

\subsection{Realización racionalizada de árbol}

La lectura electrodébil de esta coordenatización define
$c_W^2=D$, $s_W^2=1-D$ y $e_g^2=4\pi\alpha$. La raíz positiva fija
$g=e_g/s_W$ y $g'=e_g/c_W$. Se obtiene
\begin{equation}
 g^2=\frac{4\pi\alpha}{1-D},\qquad
 g'^2=\frac{4\pi\alpha}{D}.
 \label{amp-higgs-calibre}
\end{equation}
La cantidad $e_g$ es un acoplamiento adimensional en la normalización
racionalizada, distinto de una sección dimensional de carga.
La convención del potencial escalar es
\begin{equation}
 \begin{gathered}
 V(H)=-\mu_H^2H^\dagger H+\lambda_H(H^\dagger H)^2,\\
 m_H^2=2\lambda_Hv^2,\qquad m_W^2=\frac{g^2v^2}{4}.
 \end{gathered}
 \label{amp-higgs-potencial}
\end{equation}
La misma $v$ interviene en ambas relaciones. Al dividirlas, su escala se
cancela y $m_H^2/m_W^2=8\lambda_H/g^2$. Con
\eqref{amp-higgs-cocientes} resulta
\begin{equation}
 \boxed{\lambda_H=\frac{g^2}{8}\exp\!\left(6x+\frac{10}{3}y\right).}
 \label{amp-higgs-lambda}
\end{equation}
El factor $8$ procede de las dos normalizaciones cuadráticas de
\eqref{amp-higgs-potencial}; los coeficientes $6$ y $10/3$ proceden de
las firmas. La distinción conserva la operación que produce cada factor.
La composición recibe las salidas HMT previas y usa esta realización física
de árbol a escala común; un cambio de normalización del campo, de escala o
de esquema se transporta también en sus relaciones de masas y acoplamientos.

\subsection{Inversa conjunta y dominio exacto}

\begin{theorem}[Reconstrucción angular mediante calibre e Higgs]
\label{amp-higgs-inversa-teorema}
Para $\pi>0$ fijo, la aplicación de
\eqref{amp-higgs-transporte}, \eqref{amp-higgs-calibre} y
\eqref{amp-higgs-lambda}, sobre $\alpha>0$, $x>y>0$, es inyectiva y tiene
la inversa
\begin{equation}
 \begin{aligned}
 D&=\frac{g^2}{g^2+g'^2},&
 x&=-\frac12\log D,&
 \alpha&=\frac{g^2g'^2}{4\pi(g^2+g'^2)},\\
 y&=\frac3{10}\left[\log\!\left(\frac{8\lambda_H}{g^2}\right)-6x\right].
 \end{aligned}
 \label{amp-higgs-inversa}
\end{equation}
Su imagen exacta está formada por las ternas $g,g',\lambda_H>0$ tales que,
con $D$ y $x$ reconstruidos como arriba,
\begin{equation}
 \frac{g^2}{8}e^{6x}<\lambda_H<\frac{g^2}{8}e^{28x/3}.
 \label{amp-higgs-dominio}
\end{equation}
\end{theorem}
\begin{proof}
Las dos ecuaciones de calibre dan
$g^2/(g^2+g'^2)=D$ y
$g^2g'^2/(g^2+g'^2)=4\pi\alpha$. El logaritmo real recupera $x$ y
la ecuación escalar recupera $y$. Las condiciones $g,g'>0$ implican
$0<D<1$ y $x>0$. La desigualdad~\eqref{amp-higgs-dominio} equivale,
mediante el logaritmo estrictamente creciente, a
$0<\log(8\lambda_H/g^2)-6x<10x/3$, es decir, $0<y<x$.
Recíprocamente, las fórmulas reconstruidas producen
$4\pi\alpha/(1-D)=g^2$, $4\pi\alpha/D=g'^2$ y
$g^2e^{6x+10y/3}/8=\lambda_H$. Las raíces positivas recuperan los
acoplamientos originales. Esto prueba necesidad, suficiencia y unicidad.
\end{proof}

La inversión anterior describe exactamente esta aplicación entre lectores.
La sección generada por HMT conserva además la relación de
\eqref{amp-higgs-angulos}, que se expresa como
\begin{equation}
 -\frac12\log\!\left(\frac{g^2}{g^2+g'^2}\right)
 =\frac{25}{18}\frac{g^2g'^2}{g^2+g'^2}.
 \label{amp-higgs-compatibilidad-alpha}
\end{equation}
Se obtiene sustituyendo la inversa de $\alpha$ en $x=50\pi\alpha/9$.
El contraángulo conserva igualmente su ecuación de retorno anterior.
Por ello, la imagen del lector y la sección producida por el generador
son objetos distintos: las compatibilidades previas seleccionan la sección
dentro del dominio reconstruible. La operación inversa actúa después de
la generación y permite comprobar sus salidas conservadas.

\subsection{Reconstrucción Higgs--impedancia}

Para expresar el contenido orientado sin alterar los canales,
definimos
\begin{equation}
 \rho_{\rm or}=\frac{q_-}{q_+}=e^{2y},\qquad
 \Xi_H=\frac{8D^3\lambda_H}{g^2}.
 \label{amp-higgs-Xi}
\end{equation}
La ecuación~\eqref{amp-higgs-lambda} y $D^3=e^{-6x}$ dan
\begin{equation}
 \begin{aligned}
 \Xi_H&=e^{10y/3}=\rho_{\rm or}^{5/3},&
 \rho_{\rm or}&=\Xi_H^{3/5},\\
 q_-&=\sqrt D\,\Xi_H^{3/10},&
 q_+&=\sqrt D\,\Xi_H^{-3/10}.
 \end{aligned}
 \label{amp-higgs-canales}
\end{equation}
Las potencias son las raíces reales positivas. El dominio exacto
de la cámara orientada es
\begin{equation}
 0<D<1,\qquad 1<\Xi_H<D^{-5/3}.
 \label{amp-higgs-dominio-Xi}
\end{equation}
En efecto, $0<y<x$ equivale a $1<e^{10y/3}<e^{10x/3}$.
El extremo inferior hace $q_+=q_-$; el superior hace $q_-=1$.
Para las dos orientaciones juntas, $|y|<x$ equivale a
$D^{5/3}<\Xi_H<D^{-5/3}$.

Sea $F(q)=f(q^{30})=(1-q^{90})/(1-q^{120})$, la respuesta ya probada en
\eqref{iii:eq:rchannels}--\eqref{iii:eq:monotone}.
\begin{corollary}[Confluencia de las lecturas escalar y constitutiva]
\label{amp-higgs-impedancia-corolario}
En el dominio~\eqref{amp-higgs-dominio-Xi},
\begin{equation}
 \boxed{\widehat Z=
 \frac{F(\sqrt D\,\Xi_H^{3/10})}
 {F(\sqrt D\,\Xi_H^{-3/10})}.}
 \label{amp-higgs-impedancia}
\end{equation}
Las respuestas completas son
\begin{equation}
 \begin{aligned}
 r_-&=F(\sqrt D\,\Xi_H^{3/10}),&
 r_+&=F(\sqrt D\,\Xi_H^{-3/10}),\\
 \widehat\mu&=r_-^2,& \widehat\varepsilon&=r_+^2,&
 \widehat c&=(r_-r_+)^{-1}.
 \end{aligned}
 \label{amp-higgs-vacio}
\end{equation}
Además, las dos reconstrucciones satisfacen
\begin{equation}
 \frac{F^{-1}(\sqrt{\widehat\mu})}
 {F^{-1}(\sqrt{\widehat\varepsilon})}
 =\rho_{\rm or}=\Xi_H^{3/5}.
 \label{amp-higgs-confluencia}
\end{equation}
\end{corollary}
\begin{proof}
Sustituir los canales de~\eqref{amp-higgs-canales} en las cuatro lecturas
de~\eqref{iii:eq:four} prueba las primeras identidades. La inversa
constitutiva de~\eqref{iii:eq:cubic} recupera $q_-$ y $q_+$; su cociente
es $\rho_{\rm or}$. La igualdad con $\Xi_H^{3/5}$ ya está probada en
\eqref{amp-higgs-canales}.
\end{proof}

En la sección compatible con $A_{\deg}=1000\alpha$, el acoplamiento
$Q(D):=-9\log D/25$ satisface $Q(D)=4\pi\alpha$ y
$g^2=Q(D)/(1-D)$. Por tanto, la misma coordenada de Higgs se escribe
\begin{equation}
 \Xi_H=\frac{8(1-D)D^3\lambda_H}{Q(D)},\qquad
 \frac1{g^2}+\frac1{g'^2}=\frac1{Q(D)}.
 \label{amp-higgs-curva}
\end{equation}
Estas expresiones relacionan el acoplamiento total, su reparto racionalizado
y la orientación escalar en la misma sección.

\subsection{Información conservada y transporte de orientación}

La pareja de calibre recupera $\alpha$ y $x$; el cociente
$\lambda_H/g^2$ recupera después $y$. Con las marcas $P_\pm$ se reconstruyen
los operadores completos $T$ y $\mathcal V=r_+P_++r_-P_-$.
Los tres escalares determinan los autovalores; los proyectores marcados
determinan además su realización operatoria. La historia enriquecida de la
ruta sigue siendo una información anterior de mayor resolución.

Bajo $y\mapsto-y$ respecto de las marcas fijas, $D,g,g'$ permanecen
invariantes y $\Xi_H\mapsto\Xi_H^{-1}$. Se intercambian las constitutivas,
$\widehat Z\mapsto\widehat Z^{-1}$ y $\widehat c$ se conserva.
Si $\lambda_H^+$ y $\lambda_H^-$ denotan las evaluaciones en las dos
orientaciones conjugadas, se cumple
\begin{equation}
 \lambda_H^+\lambda_H^-
 =\frac{g^4}{64}D^{-6}.
 \label{amp-higgs-orientacion}
\end{equation}
La prueba consiste en multiplicar~\eqref{amp-higgs-lambda} para $y$ y $-y$;
las contribuciones orientadas se cancelan. Esta comparación conserva el
lector de firmas y cambia su argumento angular. Una conjugación simultánea
de rutas y referencias debe transportar también las firmas.

En la amplificación $T_m=T\otimes I_{9^m}$, con
$\tau_m=\operatorname{Tr}/(2\cdot9^m)$ y
$\mathcal R_m=\mathcal R\otimes I_{9^m}$, las coordenadas adecuadas son
\begin{equation}
 D=\exp\bigl(2\tau_m(\log T_m)\bigr),\qquad
 \rho_{\rm or}=\exp\bigl(-2\tau_m(\mathcal R_m\log T_m)\bigr).
 \label{amp-higgs-refinamiento}
\end{equation}
En efecto, $\log T_m=(\log T)\otimes I_{9^m}$, de modo que las dos trazas
son $-x$ y $-y$, respectivamente. La normalización elimina la multiplicidad
de la amplificación. Las lecturas de calibre, Higgs y vacío permanecen así
compatibles con estos refinamientos; el determinante ordinario ampliado
sería $D^{9^m}$ y tiene una función distinta.

La conclusión de la composición es precisa: el invariante par de calibre
y la coordenada orientada escalar reconstruyen conjuntamente la pareja
espectral que gobierna el vacío. Cuando $\lambda_H$ se evalúa mediante la
misma fórmula, la confluencia es una identidad de consistencia. Una lectura
independiente de $\lambda_H$, transportada a idénticos esquema y escala,
convierte esa igualdad en un contraste adicional. En ambos usos permanecen
explícitos las firmas, las marcas de hoja y la normalización de árbol.
